% !TEX root = ../main.tex
\section{La Búsqueda: Diseño y Sustentación de Estrategias}
\label{sec:busqueda}

En un problema de secuenciación sobre permutaciones de $M$ escenas, el tamaño del espacio combinatorio crece factorialmente como $|\mathfrak{S}_M| = M!$ (alcanzando $20! \approx 2{,}43 \times 10^{18}$ estados en ``Desenfreno de Pasiones''). Conforme a la teoría de búsqueda en $\text{COP}(\mathcal{FD})$ (Capítulo 8 de Apt) y a la literatura especializada en planificación de ensayos (Smith, 2003), la eficiencia del esquema \emph{Branch-and-Bound} (y su combinación con generación perezosa de cláusulas, \emph{Lazy Clause Generation}, en motores modernos como \texttt{Chuffed}) está gobernada por la sinergia entre el orden topológico en que se fijan las posiciones del cronograma y el criterio de selección de escenas en cada posición.

\subsection{Definición de las Ocho Estrategias de Búsqueda Exploradas}

Con el propósito de realizar una evaluación sistemática y rigurosa, ambos modelos implementan ocho estrategias de distribución parametrizadas ($E_1$--$E_8$). Sea $\sigma = (e_{(1)}, e_{(2)}, \dots, e_{(M)})$ la ordenación estática de las escenas según su \emph{peso salarial actoral ascendente} $W_e = \sum_{a \in A : p_{a,e}=1} c_a$ (desempatando por menor índice $e$), y sea $R_k = OrdenRango_k \in \{1,\dots,M\}$ el rango de la escena $Orden_k$ dentro de $\sigma$. Asimismo, defínanse el vector intercalado \emph{Outside-In} (de los extremos hacia el centro) y su homólogo reindexado por peso salarial:
\begin{align}
    \vec{z}_{\text{ext}} &= \big(Orden_1,\; Orden_M,\; Orden_2,\; Orden_{M-1},\; \dots,\; Orden_{\lceil (M+1)/2 \rceil}\big) \label{eq:vec_outside_in} \\
    \vec{r}_{\text{ext}} &= \big(R_1,\; R_M,\; R_2,\; R_{M-1},\; \dots,\; R_{\lceil (M+1)/2 \rceil}\big) \nonumber
\end{align}
Las ocho estrategias evaluadas son:

\begin{enumerate}
    \item \textbf{$E_1$ (\emph{Outside-In} por rango salarial ascendente — Estrategia principal):}
    Emplea \texttt{OrdenRangoExt} con selección \texttt{input\_order} y asignación \texttt{indomain\_min}. Instancia las posiciones alternando simétricamente entre el inicio y el final del rodaje ($1, M, 2, M-1, \dots$) y prueba en primer lugar las escenas de menor peso salarial acumulado $W_e$ en los extremos, reservando implícitamente las escenas multitudinarias y de actores costosos para el núcleo central del cronograma.

    \item \textbf{$E_2$ (\emph{Outside-In} con índice nominal de escena):}
    Emplea \texttt{OrdenExt} con \texttt{input\_order} e \texttt{indomain\_min}, conservando la topología de instanciación desde los extremos hacia el centro \eqref{eq:vec_outside_in}, pero seleccionando los valores según el identificador original $e = 1, 2, \dots, M$ sin ponderar las tarifas actorales.

    \item \textbf{$E_3$ (Secuencial cronológica de izquierda a derecha):}
    Recorre el arreglo \texttt{Orden} mediante \texttt{input\_order} y asigna con \texttt{indomain\_min}, construyendo el plan de rodaje de principio a fin ($k = 1, 2, \dots, M$) en orden numérico creciente.

    \item \textbf{$E_4$ (\emph{First-Fail} puntual sobre turnos):}
    Selecciona dinámicamente sobre \texttt{Orden} la posición $k \in K$ de menor dominio actual (\texttt{first\_fail}) y le asigna el menor índice disponible (\texttt{indomain\_min}).

    \item \textbf{$E_5$ (\emph{First-Fail} por bisección de dominio):}
    Selecciona la posición de menor dominio (\texttt{first\_fail}) y particiona su conjunto de escenas candidatas mediante bisección binaria (\texttt{indomain\_split}).

    \item \textbf{$E_6$ (Vista dual por escenas críticas centradas):}
    Invierte la perspectiva sobre el arreglo dual \texttt{PosCostoDesc} empleando orden de entrada y valor mediano (\texttt{indomain\_median}): selecciona primero las posiciones $Pos_e$ de las escenas con mayor peso salarial $W_e$ (donde actúan las estrellas más caras, como el Actor 6 con $c_6 = 40$ y el Actor 8 con $c_8 = 20$) y las ubica cerca de la mediana temporal del rodaje ($\approx M/2$).

    \item \textbf{$E_7$ (Grado ponderado por conflictos sobre \emph{Outside-In}):}
    Sobre el vector $\vec{r}_{\text{ext}}$ utiliza el selector adaptativo \texttt{dom\_w\_deg} junto con \texttt{indomain\_min}, ponderando el tamaño del dominio de cada variable entre el contador acumulado de fallos de los propagadores en que interviene.

    \item \textbf{$E_8$ (Heurística nativa \texttt{default} de cada motor):}
    Delega el control total de exploración en el buscador autónomo del solver (\emph{activity-based search} en \texttt{Gecode 6.4.0}; ramificación conflictiva VSIDS con aprendizaje de cláusulas nogood en \texttt{Chuffed 0.14.0}).
\end{enumerate}

\subsection{Sustentación Teórica de la Superioridad de la Estrategia \texorpdfstring{$E_1$}{E1}}
\label{subsec:sustento_e1}

La drástica ventaja empírica de $E_1$ (y su complemento adaptativo con aprendizaje de cláusulas en \texttt{Chuffed}) frente a las heurísticas cronológicas o de fallo temprano convencionales se sustenta en tres propiedades estructurales exactas del problema de ensayos:

\begin{itemize}
    \item \textbf{Cierre bilateral temprano de intervalos actorales bajo ramificación \emph{Outside-In}:}
    Al construir el cronograma de izquierda a derecha ($E_3$), se conoce con certeza el instante de llegada $ini_a$ de los actores que ya debutaron en $\{1,\dots,r\}$, pero su salida $fin_a$ permanece indeterminada hasta fijar las últimas posiciones cercanas a $M$, pues cualquier escena pendiente de $E_a$ aún podría ir en $r+1$ o en $M$. En consecuencia, $ultima_a$ conserva cota inferior pequeña y ningún literal de intermediación de la restricción redundante $R_1$ puede fijarse a $1$ en la mitad superior del árbol.
    En cambio, al alternar los extremos $1, M, 2, M-1, \dots$ ($E_1$ y $E_2$), a profundidad $2r$ quedan fijados simultáneamente el prefijo $\{1,\dots,r\}$ y el sufijo $\{M-r+1,\dots,M\}$, confinando las escenas restantes al bloque central $\{r+1, \dots, M-r\}$. Todo actor $a$ presente al menos una vez en el prefijo y una vez en el sufijo tiene ya cerrados ambos extremos ($primera_a \le r$ y $ultima_a \ge M-r+1$), por lo que $R_1$ deduce en tiempo constante que dicho actor permanece esperando durante \emph{todas} las escenas ajenas restantes del bloque central, elevando de inmediato $\underline{Costo}$.

    \item \textbf{Activación en profundidad $2$ de la ruptura de simetría $S_1$:}
    Mientras que en $E_3$ la desigualdad $Orden_{q+1} < Orden_M$ solo puede evaluarse por completo al llegar a la hoja de profundidad $M$ (cuando se instancia $Orden_M$), en $E_1$ las variables $Orden_{q+1}$ y $Orden_M$ son las dos primeras variables activas en instanciarse, podando exactamente el $50\%$ del árbol combinatorio en el segundo nivel de profundidad.

    \item \textbf{Alineación económica de la heurística de valor por peso salarial ascendente:}
    En toda solución óptima de baja espera, las posiciones periféricas ($1, M, 2, M-1$) deben ocuparse con escenas de pocos actores y bajo costo salarial agregado (tales como la escena $4$ con solo el Actor 1, la escena $19$ con solo el Actor 5 o la escena $16$ con solo el Actor 6), mientras que las escenas densas donde confluyen cinco o seis actores simultáneamente (como las escenas $7$ y $8$ en ``Desenfreno de Pasiones'', con peso salarial conjunto $38$ y $48$ por unidad de tiempo) deben situarse en el centro del rodaje donde los intervalos de permanencia del elenco se cruzan naturalmente. Al ordenar los valores por $W_e$ creciente, $E_1$ descubre en las primeras ramas una solución incumbente de costo muy próximo al óptimo global ($871$ en Parte 1 y $880$ en Parte 2), maximizando el poder de poda de la cota superior $Costo < Costo_{\text{inc}}$ durante el resto de la demostración de optimalidad.
\end{itemize}
